\documentclass[11pt,a4paper]{article}
\usepackage[margin=25mm]{geometry}
\usepackage[T1]{fontenc}
\usepackage[utf8]{inputenc}
\usepackage{lmodern,amsmath,amssymb,booktabs,microtype}
\usepackage[hidelinks]{hyperref}
\newcommand{\dd}{\mathrm d}
\newcommand{\ii}{\mathrm i}
\newcommand{\Tr}{\operatorname{Tr}}
\newcommand{\R}{\mathbb R}
\title{Quantizing the sextic Q-ball sector:\\charge sectors, collective phase and a lattice matching prescription}
\date{5 October 2026 --- English working draft v0.1}
\begin{document}
\maketitle
\begin{abstract}
We formulate a canonical quantum version of the classical complex scalar
model with dimensionless potential $U(s)=s-s^2+s^3/2$.
On a fixed finite spatial lattice, a positive mass matrix and a nonnegative
stiffness matrix define a Hamiltonian with exact integer global charge sectors.
An explicit action normalization $\zeta$ relates the classical charge to
integer quantum charge through $q=\zeta n$; the leading semiclassical energy
is $E_{\mathrm{cl}}(\zeta n)/\zeta$ rather than an unqualified evaluation at $n$.
Quantization of the collective phase reproduces this relation. We give the
quadratic rotating-frame Hamiltonian, identify its charge and symmetry
constraints, and specify a correlator comparison for the separate numerical
quantization effort. Spatial dimension changes the field theory's power
counting, profiles and renormalization requirements. This draft establishes
a regulated quantum formulation and a matching prescription; it reports no
new quantum spectrum, binding energy, loop mass or experimental particle identification.
\end{abstract}

\noindent\textbf{Provenance and status.}
This is a new English derivation, not a translation of the existing Q-ball
ladder paper or the coordinator's broader particle-model manuscript.
The potential and charge sign follow the current project convention.
The separate Claude task QUANT-1 proposes Euclidean correlator measurements;
its results are not assumed or reproduced here. Attribution remains

\section{What is being quantized}
We first quantize the scalar field on a fixed background. The dimensionless
coordinates are $\tau=m t$, $y=m x$, with $\hbar=c=1$ unless stated otherwise.
Let $d$ denote spatial dimension, and define
\begin{equation}\label{eq:action}
 S=\frac1\zeta\int\dd\tau\,\dd^dy\,
 \left[|\partial_\tau\phi|^2-|\nabla\phi|^2-U(|\phi|^2)\right],
 \qquad U(s)=s-s^2+\tfrac12s^3,\quad \zeta>0.
\end{equation}
Here $\zeta$ is part of the quantum model. Multiplying a classical action
by a constant leaves its classical equations unchanged, but changes the
weight $\exp(\ii S)$ and the size of quantum fluctuations. Taking
$\zeta=1$ defines a possible bare pilot; it is not a weak-coupling argument.
No value of $m$ in physical units or phenomenological value of $\zeta$ has
been determined by the present classical project.

For a canonically normalized dimensionful field,
\begin{equation}\label{eq:rescale}
 \Phi=\frac{m^{(d-1)/2}}{\sqrt\zeta}\phi,
 \quad V(\Phi)=m^2|\Phi|^2-\lambda_4|\Phi|^4+\lambda_6|\Phi|^6,
 \quad \lambda_4=\zeta m^{3-d},\quad
 \lambda_6=\tfrac12\zeta^2m^{4-2d}.
\end{equation}
These are bare, tree-level relations at the specified regulator. Holding
renormalized parameters fixed while changing the regulator requires matching.
The equivalent loss of action normalization under classical field rescaling
is explicitly discussed for a sextic model in Ref.~\cite{xie}, Sec.~2.5,
Eqs.~(30)--(31). Its charge sign and radial amplitude conventions differ
from ours; numerical charges are not imported.

This step introduces standard quantum mechanics through the Hilbert space
and commutation relations below. It does not derive quantum mechanics from
classical geometry. The dynamical Regge and Maxwell sectors of the separate
network proposal are not included in this scalar Hamiltonian.

\section{An explicit regulated quantum Hamiltonian}
Write $\phi_i=(q_{1i}+\ii q_{2i})/\sqrt2$ on $N_s$ spatial sites. Choose
positive site volumes $w_i$, $W=\operatorname{diag}(w_i)$, and a real symmetric
nonnegative stiffness matrix $K$ representing the gradient quadratic form.
Boundary conditions and the geometry are fixed. Then
\begin{align}
 L&=\frac1\zeta\left[\frac12\sum_{a=1}^2\dot q_a^T W\dot q_a
 -\frac12\sum_{a=1}^2q_a^TKq_a-\sum_iw_iU(s_i)\right],
 &s_i&=\frac{q_{1i}^2+q_{2i}^2}{2},\\
 p_a&=\zeta^{-1}W\dot q_a,\qquad
 [\hat q_{ai},\hat p_{bj}]=\ii\delta_{ab}\delta_{ij},\qquad
 \mathcal H=L^2(\R^{2N_s}),\\
 \hat H/m&=\frac\zeta2\sum_{a=1}^2\hat p_a^TW^{-1}\hat p_a
 +\frac1{2\zeta}\sum_{a=1}^2\hat q_a^TK\hat q_a
 +\frac1\zeta\sum_iw_iU(\hat s_i).\label{eq:ham}
\end{align}
The coordinates commute, so $U(\hat s_i)$ has no ordering ambiguity in this
Cartesian representation. The identity
$U(s)=s[(s-1)^2+1]/2\geq0$ and the positive kinetic terms give a lower-bounded
finite-lattice Hamiltonian. This is a specified quantum system even before
its eigenvalues are computed. Continuum renormalization is a further problem.
A consistent non-diagonal positive mass matrix could replace $W$ with the
corresponding inverse in the kinetic term; a proposed irregular lattice must
supply its actual matrices, not reuse cubic-grid weights without checking.

Global phase rotations have the generator
\begin{equation}\label{eq:charge}
 \hat N=\sum_i(\hat q_{1i}\hat p_{2i}-\hat q_{2i}\hat p_{1i}),
 \qquad [\hat H,\hat N]=0,\qquad [\hat N,\hat\phi_i]=\hat\phi_i.
\end{equation}
Single-valued wavefunctions under a $2\pi$ global rotation give
$n\in\mathbb Z$. These are charges, not a fixed number of constituents:
particle--antiparticle pairs can occur within one charge sector.
The exact projector and sector partition function are
\begin{equation}\label{eq:projector}
 P_n=\frac1{2\pi}\int_0^{2\pi}\dd\theta\,
 e^{\ii\theta(\hat N-n)},\qquad
 Z_n(T)=\Tr(P_n e^{-T\hat H/m}).
\end{equation}
The observable sector energy is measured relative to the vacuum:
\begin{equation}
 M_n/m=-\lim_{T\to\infty}T^{-1}\log[Z_n(T)/Z_0(T)].
\end{equation}
This formula is a definition, not a proposed efficient Monte Carlo estimator.

\section{Integer charge and the collective phase}
For $\phi=e^{\ii\omega\tau}f(y)$ with real localized $f$,
\begin{equation}
 q=2\omega I,\quad I=\int f^2\dd^dy,\quad
 E_{\mathrm{cl}}=\omega^2I+\int[|\nabla f|^2+U(f^2)]\dd^dy.
\end{equation}
The classical solution's physical charge and energy are
\begin{equation}\label{eq:matching}
 n=\frac q\zeta,\qquad \frac{E_{\mathrm{physical}}}m=
 \frac{E_{\mathrm{cl}}}\zeta.
\end{equation}
To derive the same result collectively, freeze a profile temporarily and
allow $\phi=e^{\ii\alpha(\tau)}f$. If
$\mathcal V[f]=\int[|\nabla f|^2+U(f^2)]\dd^dy$, then
\begin{equation}
 L_\alpha=\zeta^{-1}(I\dot\alpha^2-\mathcal V),\quad
 p_\alpha=2I\dot\alpha/\zeta,\quad
 \hat p_\alpha=-\ii\partial_\alpha,\quad
 \psi_n(\alpha)=\frac{e^{\ii n\alpha}}{\sqrt{2\pi}}.
\end{equation}
Consequently the frozen-profile rotor energy is
\begin{equation}
 \frac{H_\alpha}m=\frac{\zeta n^2}{4I}+\frac{\mathcal V}\zeta.
\end{equation}
Minimizing over profiles at fixed $n$ is precisely the classical fixed-charge
variational problem at $q=\zeta n$. If the chosen classical solution is the
relevant lowest-energy configuration in that charge sector and the
semiclassical expansion is controlled, this gives
\begin{equation}\label{eq:mass}
 \frac{M_n}m=\frac1\zeta E_{\mathrm{cl}}(\zeta n)
 +\Delta_1(\zeta n;\Lambda,\text{matching})+O(\zeta),
\end{equation}
For a merely local or metastable branch, the right-hand side instead
describes that branch's semiclassical energy; it is not automatically the
lowest sector energy $M_n$ defined by Eq.~\eqref{eq:projector}. Competing
single-object and multi-object configurations must be compared. In that
case the spacing formula below also applies only to the selected branch.
Here $\Delta_1$ denotes a renormalized vacuum-subtracted fluctuation
correction, not yet evaluated. This expansion is at fixed classical
$q=\zeta n$ as $\zeta\to0$, hence $n\to\infty$.
It is not a controlled expansion at $n=2$ or $3$ simply because a classical
profile exists at another charge.

In a smooth semiclassical branch with $\dd E_{\mathrm{cl}}/\dd q=\omega$,
\begin{equation}\label{eq:spacing}
 \frac{M_{n+1}^{(0)}-M_n^{(0)}}m
 =\omega+\frac\zeta2\frac{\dd\omega}{\dd q}+O(\zeta^2)
 =\omega+\frac\zeta{2q'(\omega)}+O(\zeta^2).
\end{equation}
This expansion fails near a turning point where $q'(\omega)=0$ and cannot
replace a finite difference across a branch change. Negative curvature
between charge sectors is not by itself a negative oscillation mode within
one fixed-charge sector. The expression concerns leading classical masses;
loop corrections change the spacing at the next relevant order.

An exact charge eigenstate has $\langle\hat\phi\rangle=0$ by the charge
selection rule. The familiar rotating classical mean field is therefore a
semiclassical description with a phase reference, not literally the one-point
function of an exact fixed-charge eigenstate. Its density and correlation
functions can still encode the localized object. In a translation-invariant
system a momentum eigenstate also has a delocalized center of mass; a localized
Q-ball is described by a collective-coordinate wavepacket or conditional
internal observables. On a finite lattice translations have only the lattice
symmetry and continuum center-of-mass formulas need a convergence check.

\section{Quantizing internal fluctuations without double counting}
In the continuum classical background convention, put
\begin{equation}
 \phi=e^{\ii\omega\tau}\left[f+\sqrt\zeta\,(u+\ii v)/\sqrt2\right],
 \quad L_-=-\Delta+1-2f^2+\tfrac32f^4-\omega^2,
 \quad L_+=-\Delta+1-6f^2+\tfrac{15}{2}f^4-\omega^2.
\end{equation}
The quadratic rotating-frame Lagrangian and momenta are
\begin{align}
 L^{(2)}&=\int\dd^dy\left[\tfrac12(\dot u^2+\dot v^2)
 +\omega(u\dot v-v\dot u)-\tfrac12uL_+u-\tfrac12vL_-v\right],\\
 p_u&=\dot u-\omega v,\qquad p_v=\dot v+\omega u.
\end{align}
Its Hamiltonian is the quadratic part of $H/m-\omega N$:
\begin{equation}\label{eq:k2}
 K^{(2)}=\frac12\int\dd^dy\left[(p_u+\omega v)^2+(p_v-\omega u)^2
 +uL_+u+vL_-v\right].
\end{equation}
It generates $\ddot u-2\omega\dot v+L_+u=0$ and
$\ddot v+2\omega\dot u+L_-v=0$. Diagonalizing only $L_+$ or adding
independent sideband frequencies misses the canonical mixing.
A mode basis must preserve the symplectic form
$\int(\delta u\wedge\delta p_u+\delta v\wedge\delta p_v)$.

The phase zero mode and the $d$ translation modes are collective coordinates,
not ordinary positive-frequency oscillators. The linear charge constraint is
\begin{equation}
 \delta N=\sqrt{2/\zeta}\int f(2\omega u+\dot v)\dd^dy=0.
\end{equation}
Charge must also be treated at the order needed for the energy calculation;
a linear constraint alone is not a complete one-loop fixed-charge prescription.
A negative or growing physical mode invalidates a stable harmonic-vacuum
interpretation for that background. A classical embedded bound mode must be
normalizable and correctly normalized, but its linear trapping does not by
itself suppress nonlinear or quantum decay channels.

For a background that passes the constrained stability checks, the positive
canonical modes can be promoted to creation and annihilation operators.
The resulting zero-point sums require continuum scattering contributions,
vacuum subtraction and counterterms in the same regulator and matching
scheme. No number for $\Delta_1$ follows merely from the existing BIC frequency.
Ref.~\cite{graham} develops quantum corrections for a different potential;
its reported correction is not imported into this sextic model.

\section{A matching prescription for QUANT-1}
Euclidean zero-density sampling and charged correlators can determine sector
energies without inserting a real chemical potential into the sampling
weight. This is spectroscopy, not a simulation of real-time collisions.
Choose an operator $O_n$ with $[N,O_n]=nO_n$, and a basis of local and
spatially extended operators if possible. At zero temperature,
\begin{equation}
 C_n(t)=\langle0|O_n^\dagger(t)O_n(0)|0\rangle
 =\sum_a|\langle n,a|O_n|0\rangle|^2
 e^{-[E_{n,a}-E_0]t/m}.
\end{equation}
A missing plateau can mean poor overlap or excited-state contamination;
it does not demonstrate the absence of a bound state.

The following information is part of the comparison, not optional metadata:
$\zeta$, bare versus renormalized mass and couplings, spatial geometry,
spacings and volumes, time extent, operator definitions, autocorrelation,
uncertainty estimates and the fitting window. For the classical comparison
in Eq.~\eqref{eq:mass}, use the same matched theory and $q=\zeta n$.
The physical one-particle mass $M_1$ must be measured rather than silently
identified with the bare coefficient $m$.

An exact free-field check is particularly useful. If $\Phi_s$ is a fixed
linear spatial smearing of a free complex Gaussian field and
$O_n=\Phi_s^n/\sqrt{n!}$, Wick contraction gives
\begin{equation}\label{eq:wick}
 C_n(t)=G_s(t)^n,\qquad
 G_s(t)=\langle\Phi_s^\dagger(t)\Phi_s(0)\rangle.
\end{equation}
At finite periodic Euclidean time, even a single-mode $G_s$ contains forward
and backward propagation. Its $n$th power is generally not a single cosh
with mass $nM_1$. Mixed thermal terms can mimic anomalously low energies;
the free check must use the actual time extent and operator.

Define $B_2=2M_1-M_2$ and $B_3=3M_1-M_3$.
Positive values with controlled uncertainty and spatial-volume behavior are
necessary comparisons against free constituents. A three-charge bound state
must also lie below $M_2+M_1$ if a two-charge bound state exists; larger charges
require all relevant partitions. Finite-volume interaction shifts of scattering
states can lie below the free threshold without an infinite-volume bound pole.
A binding claim therefore needs volume dependence as well as a plateau.
We do not set a numerical binding tolerance before the estimator and budget
are specified. QUANT-1 remains a separate numerical task, not duplicated here.

\section{What changes with spatial dimension}
Canonical charge quantization and Eqs.~\eqref{eq:ham}--\eqref{eq:matching}
apply for any integer $d$ at a fixed finite regulator. Profiles and integration
measures do not: the continuum radial measure is
$S_{d-1}r^{d-1}\dd r$, and the angular spectrum is that of $S^{d-1}$.
The one-dimensional problem on the whole line requires its own parity
sectors rather than importing three-dimensional angular labels.

\begin{center}
\begin{tabular}{@{}llll@{}}
\toprule
Spatial $d$ & $[\Phi]$ & $[\lambda_4]$ & $[\lambda_6]$\\
\midrule
1 & 0 & 2 & 2\\
2 & $1/2$ & 1 & 0\\
3 & 1 & 0 & $-2$\\
$d$ & $(d-1)/2$ & $3-d$ & $4-2d$\\
\bottomrule
\end{tabular}
\end{center}
Brackets denote mass dimensions. In $2+1$ dimensions the sextic coupling is
marginal by power counting; in $3+1$ it is irrelevant and the polynomial is
naturally a cutoff effective field theory. A cutoff-independent ultraviolet
completion is not specified by the classical sextic coefficients. Additional
operators and renormalization conditions must be included to the desired EFT
accuracy. Lower-dimensional counterterms and stability observations cannot be
transferred unchanged to $3+1$ dimensions.

A radial three-dimensional numerical profile is not a one-dimensional quantum
field theory: discarding angular quantum fluctuations defines a truncation.
For an additional compact spatial coordinate the transverse modes, radius,
normalization and thresholds must be retained. An internal component index is
not automatically that coordinate. None of these variants derives spatial
dimension from energy or generates the time variable used for quantization.

\section{First result and next calculation}
The completed first step is a concrete, lower-bounded finite-lattice quantum
Hamiltonian, integer charge projection, and a charge-normalized classical
matching formula. The rotating-frame fluctuation Hamiltonian identifies the
next semiclassical calculation without confusing zero modes and oscillators.
This is a quantum formulation of the scalar sector, not yet a solved quantum
Q-ball spectrum or a quantization of the full evolving network.

The next numerical result should come from the already assigned QUANT-1
spectroscopy with its action normalization and free correlator controls fixed.
A complementary later semiclassical calculation can use one identified stable
large-charge profile, its complete constrained fluctuation problem and a
specified renormalization prescription. It should not extrapolate a few
classical mode frequencies to a small-charge mass spectrum.
Elementary excitations of this scalar quantization are bosonic; spin one-half,
color confinement and Standard Model identities require additional structure
and are not outcomes of integer Q-ball charge.

\begin{thebibliography}{9}
\bibitem{xie} Q.-X. Xie, P. M. Saffin, A. Tranberg and S.-Y. Zhou,
\emph{Quantum corrected Q-ball dynamics}, arXiv:2312.01139v1 (2023),
\href{https://arxiv.org/html/2312.01139v1}{primary text}.
Targeted reading: Secs.~2.1, 2.4 and 2.5, especially Eqs.~(1)--(8),
(26)--(31), (35)--(36). The reported real-time scheme is an inhomogeneous
Hartree approximation; its numerical results are not reproduced here.
\bibitem{graham} N. Graham, \emph{Quantum Corrections to Q-Balls},
arXiv:hep-th/0105009v2 (2001),
\href{https://arxiv.org/html/hep-th/0105009v2}{primary text}.
Targeted reading: abstract and Sec.~I through the model definition,
Eqs.~(1)--(2). The cubic--quartic potential differs from ours; no full
reassessment of its one-loop calculation is claimed.
\end{thebibliography}
\end{document}
